\section{2021级工科数学分析（二）期末考试 A 卷}

\subsection{试题}

\paragraph{一、叙述定义（共 5 小题，每小题 2 分，共 10 分）}
\begin{enumerate}
  \item 二元函数的偏导数。
  \item 高斯公式。
  \item 条件收敛级数。
  \item 函数列的一致收敛。
  \item 二阶线性齐次常微分方程 $L[y]=0$ 的基础解系。
\end{enumerate}

\paragraph{二、计算题（共 4 小题，每小题 10 分，共 40 分）}
\begin{enumerate}
  \item 设 $u=z^2+2z$，且 $z=z(x,y)$ 由方程 $xe^x-ye^y=ze^z$（$z\ne 1$）所确定，求 $du$ 和 $\dfrac{\partial^2u}{\partial x^2}$。
  \item 计算
  \[
    \int_1^2 dx\int_{\sqrt{x}}^x \sin\frac{\pi x}{2y}\,dy
    +\int_2^4 dx\int_{\sqrt{x}}^2 \sin\frac{\pi x}{2y}\,dy .
  \]
  \item 计算 $\iint_\Sigma x\,dy\,dz$，其中 $\Sigma$ 为 $x^2+y^2=1$，$z=x+2$ 和 $z=0$ 所围立体的表面，取外侧。
  \item 计算
  \[
    \int_L e^{\sqrt{x^2+y^2}}\,ds,
  \]
  其中 $L$ 为圆周 $x^2+y^2=a^2$（$a>0$）、直线 $y=x$ 与 $x$ 轴在第一象限内所围成的扇形的整个边界。
\end{enumerate}

\paragraph{三、解答题（共 3 小题，每小题 10 分，共 30 分）}
\begin{enumerate}
  \item 设曲线积分
  \[
    \int_L xy^2\,dx+y\varphi(x)\,dy
  \]
  与路径无关，其中 $\varphi(x)$ 具有连续的导数，且 $\varphi(0)=0$，计算
  \[
    \int_{(0,0)}^{(1,1)} xy^2\,dx+y\varphi(x)\,dy .
  \]
  \item 求级数 $\displaystyle \sum_{n=1}^\infty \frac{n(n+1)}{2^n}$ 的和。
  \item 求方程 $y''+4y'+4y=\cos 2x$ 的通解。
\end{enumerate}

\paragraph{四、应用题（10 分）}
造一容积为 $V$ 的长方形无盖铝盆，怎样设计尺寸，使它的表面积最小？

\paragraph{五、证明题（10 分）}
证明级数 $\displaystyle \sum_{n=1}^\infty x^n(1-x)^2$ 在 $[0,1]$ 上一致收敛。

\subsection{答案与解析}

\paragraph{一、叙述定义}
\begin{enumerate}
  \item 若极限
  \[
    \lim_{\Delta x\to 0}\frac{f(x_0+\Delta x,y_0)-f(x_0,y_0)}{\Delta x}
  \]
  存在，则称其为 $f$ 在 $(x_0,y_0)$ 处关于 $x$ 的偏导数，记为 $f_x(x_0,y_0)$；关于 $y$ 的偏导数类似。
  \item 若闭区域 $V$ 的边界为分片光滑闭曲面 $\Sigma$，取外侧，且 $P,Q,R$ 在 $V$ 上有连续一阶偏导数，则
  \[
    \iiint_V\left(\frac{\partial P}{\partial x}+\frac{\partial Q}{\partial y}+\frac{\partial R}{\partial z}\right)\,dV
    =
    \iint_\Sigma P\,dy\,dz+Q\,dz\,dx+R\,dx\,dy .
  \]
  \item 若 $\sum a_n$ 收敛而 $\sum |a_n|$ 发散，则称 $\sum a_n$ 条件收敛。
  \item 若函数列 $\{f_n(x)\}$ 在 $I$ 上逐点收敛于 $f(x)$，且对任意 $\varepsilon>0$，存在 $N$，使当 $n>N$ 时，对一切 $x\in I$ 都有
  \[
    |f_n(x)-f(x)|<\varepsilon,
  \]
  则称 $\{f_n\}$ 在 $I$ 上一致收敛于 $f$。
  \item 若二阶线性齐次方程 $L[y]=0$ 的两个解 $\varphi_1,\varphi_2$ 线性无关，则 $\{\varphi_1,\varphi_2\}$ 称为该方程的一个基础解系。
\end{enumerate}

\paragraph{二、计算题}
\begin{enumerate}
  \item 令 $F(x,y,z)=xe^x-ye^y-ze^z$。由隐函数求导，
  \[
    z_x=\frac{x+1}{z+1}e^{x-z},\qquad
    z_y=-\frac{y+1}{z+1}e^{y-z}.
  \]
  题面括号中写作 $z\ne 1$，而隐函数求导实际要求 $F_z=-(z+1)e^z\ne0$，即 $z\ne-1$；以下按该条件计算。
  因为 $u=z^2+2z$，所以
  \[
    du=2(x+1)e^{x-z}\,dx-2(y+1)e^{y-z}\,dy,
  \]
  \[
    \frac{\partial^2u}{\partial x^2}
    =
    \frac{2e^{x-z}\bigl[(z+1)(x+2)-(x+1)^2e^{x-z}\bigr]}{z+1}.
  \]
  \item 交换积分次序，积分区域可写为 $1\le y\le 2,\ y\le x\le y^2$，故
  \[
    \int_1^2 dy\int_y^{y^2}\sin\frac{\pi x}{2y}\,dx
    =-\frac{2}{\pi}\int_1^2 y\cos\frac{\pi y}{2}\,dy
    =\frac{4(\pi+2)}{\pi^3}.
  \]
  \item 由高斯公式，取 $P=x,Q=R=0$，则
  \[
    \iint_\Sigma x\,dy\,dz=\iiint_\Omega 1\,dV
    =\iint_{x^2+y^2\le 1}(x+2)\,dx\,dy=2\pi.
  \]
  原答案抽取文本末行疑似将结果误写为 $4\pi$，本稿按高斯公式校为 $2\pi$。
  \item 边界由三段组成：$L_1:y=0,0\le x\le a$；$L_2:y=x,0\le x\le a/\sqrt2$；$L_3:x=a\cos\theta,y=a\sin\theta,0\le\theta\le\pi/4$。于是
  \[
    \int_{L_1}e^{\sqrt{x^2+y^2}}\,ds=e^a-1,\quad
    \int_{L_2}e^{\sqrt{x^2+y^2}}\,ds=e^a-1,
  \]
  \[
    \int_{L_3}e^{\sqrt{x^2+y^2}}\,ds=\frac{\pi a}{4}e^a.
  \]
  因而
  \[
    \int_L e^{\sqrt{x^2+y^2}}\,ds=2(e^a-1)+\frac{\pi a}{4}e^a.
  \]
\end{enumerate}

\paragraph{三、解答题}
\begin{enumerate}
  \item 设 $P=xy^2,Q=y\varphi(x)$。路径无关给出
  \[
    P_y=Q_x,\qquad 2xy=y\varphi'(x),
  \]
  从而 $\varphi'(x)=2x$，又 $\varphi(0)=0$，故 $\varphi(x)=x^2$。沿折线 $(0,0)\to(1,0)\to(1,1)$ 计算，
  \[
    \int_{(0,0)}^{(1,1)} xy^2\,dx+y\varphi(x)\,dy
    =\int_0^1 y\,dy=\frac12.
  \]
  \item 对 $|x|<1$，设
  \[
    S(x)=\sum_{n=1}^\infty n(n+1)x^n.
  \]
  由 $\sum_{n=1}^\infty nx^{n-1}=1/(1-x)^2$ 可得
  \[
    S(x)=\frac{2x}{(1-x)^3}.
  \]
  取 $x=1/2$ 得
  \[
    \sum_{n=1}^\infty \frac{n(n+1)}{2^n}=S\left(\frac12\right)=8.
  \]
  \item 齐次方程特征根为 $r=-2$（二重根），故
  \[
    y_h=(C_1+C_2x)e^{-2x}.
  \]
  设特解 $y_p=A\sin 2x+B\cos 2x$，代入得 $A=1/8,B=0$，所以
  \[
    y=(C_1+C_2x)e^{-2x}+\frac18\sin 2x.
  \]
\end{enumerate}

\paragraph{四、应用题}
设长方形无盖铝盆底边为 $x,y$，高为 $h$。约束为 $xyh=V$，表面积
\[
  S=xy+2xh+2yh=xy+\frac{2V}{y}+\frac{2V}{x}.
\]
由对称性和驻点方程可得 $x=y$，令 $x=y=t$，则
\[
  S(t)=t^2+\frac{4V}{t},\qquad S'(t)=2t-\frac{4V}{t^2}=0.
\]
故 $t^3=2V$，$h=V/t^2=t/2$。即底面为边长 $\sqrt[3]{2V}$ 的正方形，高为 $\sqrt[3]{2V}/2$ 时表面积最小。

\paragraph{五、证明题}
记 $u_n(x)=x^n(1-x)^2$。在 $[0,1]$ 上，
\[
  u_n'(x)=nx^{n-1}(1-x)^2-2x^n(1-x),
\]
其最大值在 $x=n/(n+2)$ 处取得，且
\[
  0\le u_n(x)\le \left(\frac{n}{n+2}\right)^n\left(\frac{2}{n+2}\right)^2\le \frac{4}{n^2}.
\]
由于 $\sum 4/n^2$ 收敛，由 Weierstrass 判别法知
\[
  \sum_{n=1}^\infty x^n(1-x)^2
\]
在 $[0,1]$ 上一致收敛。

\section{2021级工科数学分析（二）期末考试 B 卷}

\subsection{试题}

\paragraph{一、叙述定义（共 5 小题，每小题 2 分，共 10 分）}
\begin{enumerate}
  \item 二元函数的方向导数。
  \item 斯托克斯公式。
  \item 绝对收敛级数。
  \item 函数项级数的一致收敛。
  \item $n$ 阶线性齐次常微分方程 $L[y]=0$ 的基础解系。
\end{enumerate}

\paragraph{二、计算题（共 4 小题，每小题 10 分，共 40 分）}
\begin{enumerate}
  \item 设 $z=\ln(u^2+v^2)$，而 $u=e^{x^2+y^2}$，$v=x^2+y$，求 $\dfrac{\partial z}{\partial x}$ 和 $\dfrac{\partial^2z}{\partial x\partial y}$。
  \item 计算
  \[
    \int_{-4}^0 dx\int_0^{2x+8}(8-y)^{1/3}\,dy
    +\int_0^4 dx\int_0^{-2x+8}(8-y)^{1/3}\,dy .
  \]
  \item 计算
  \[
    \iint_\Sigma (z^2+x)\,dy\,dz-z\,dx\,dy,
  \]
  其中 $\Sigma$ 为旋转抛物面 $z=\frac12(x^2+y^2)$ 介于平面 $z=0$ 与 $z=2$ 之间的部分的下侧。
  \item 计算
  \[
    I=\int_C (z-y)\,dx+(x-z)\,dy+(x-y)\,dz,
  \]
  其中 $C$ 是曲线
  \[
    \begin{cases}
      x^2+y^2=1,\\
      x-y+z=2,
    \end{cases}
  \]
  从 $z$ 轴正向往 $z$ 轴负向看 $C$ 的方向是顺时针的。
\end{enumerate}

\paragraph{三、解答题（共 3 小题，每小题 10 分，共 30 分）}
\begin{enumerate}
  \item 设曲线积分
  \[
    \int_L x\varphi(y)\,dx+x^2y\,dy
  \]
  与路径无关，其中 $\varphi(y)$ 具有连续的导数，且 $\varphi(0)=1$，计算
  \[
    \int_{(0,0)}^{(1,1)}x\varphi(y)\,dx+x^2y\,dy.
  \]
  \item 求级数
  \[
    \sum_{n=2}^\infty (-1)^n\frac{x^n}{n(n-1)}
  \]
  的收敛域，并求和函数。
  \item 求方程 $y''+y'-2y=-2\sin x$ 的通解。
\end{enumerate}

\paragraph{四、应用题（10 分）}
在椭圆 $x^2+4y^2=4$ 上求一点，使其到直线 $2x+3y-6=0$ 的距离最短。

\paragraph{五、证明题（10 分）}
证明级数 $\displaystyle \sum_{n=1}^\infty x^2e^{-nx}$ 在 $(0,+\infty)$ 内一致收敛。

\subsection{答案与解析}

\paragraph{一、叙述定义}
\begin{enumerate}
  \item 若极限
  \[
    \lim_{\rho\to0^+}\frac{f(x_0+\rho\cos\alpha,y_0+\rho\sin\alpha)-f(x_0,y_0)}{\rho}
  \]
  存在，则称其为 $f$ 在 $(x_0,y_0)$ 处沿方向 $\vec l=(\cos\alpha,\sin\alpha)$ 的方向导数。
  \item 若 $\Sigma$ 是分片光滑有向曲面，边界为正向闭曲线 $L$，且 $P,Q,R$ 有连续一阶偏导数，则
  \[
    \oint_L P\,dx+Q\,dy+R\,dz
    =
    \iint_\Sigma
    \begin{vmatrix}
      dy\,dz & dz\,dx & dx\,dy\\
      \partial_x & \partial_y & \partial_z\\
      P & Q & R
    \end{vmatrix}.
  \]
  \item 若级数 $\sum a_n$ 满足 $\sum |a_n|$ 收敛，则称 $\sum a_n$ 绝对收敛。
  \item 若函数项级数 $\sum u_n(x)$ 的部分和 $S_n(x)$ 在区间 $I$ 上一致收敛于 $S(x)$，即对任意 $\varepsilon>0$，存在 $N$，使 $n>N$ 时对一切 $x\in I$ 有
  \[
    |S_n(x)-S(x)|<\varepsilon,
  \]
  则称该级数在 $I$ 上一致收敛。
  \item $n$ 阶线性齐次常微分方程 $L[y]=0$ 的 $n$ 个线性无关解 $y_1,\ldots,y_n$ 构成的解组称为一个基础解系。
\end{enumerate}

\paragraph{二、计算题}
\begin{enumerate}
  \item 令
  \[
    u=e^{x^2+y^2},\qquad v=x^2+y,\qquad A=u^2+v^2.
  \]
  则
  \[
    z_x=\frac{4x(u^2+v)}{A},
  \]
  且
  \[
    z_{xy}=
    \frac{4x\bigl[(4yu^2+1)A-(u^2+v)(4yu^2+2v)\bigr]}{A^2}.
  \]
  \item 积分区域可改写为
  \[
    0\le y\le8,\qquad \frac{y-8}{2}\le x\le\frac{8-y}{2}.
  \]
  因此原式为
  \[
    \int_0^8(8-y)^{1/3}(8-y)\,dy
    =\int_0^8(8-y)^{4/3}\,dy
    =\frac{384}{7}.
  \]
  \item 对图形 $z=\frac12(x^2+y^2)$ 取下侧，其向量面积元为 $(x,y,-1)\,dx\,dy$，投影区域为 $x^2+y^2\le4$。故
  \[
    \iint_\Sigma (z^2+x)\,dy\,dz-z\,dx\,dy
    =
    \iint_{x^2+y^2\le4}(xz^2+x^2+z)\,dx\,dy.
  \]
  奇函数项积分为 $0$，且 $z=(x^2+y^2)/2$，所以
  \[
    I=\iint_{x^2+y^2\le4}\left(x^2+\frac{x^2+y^2}{2}\right)\,dx\,dy=8\pi.
  \]
  \item 令 $\vec F=(z-y,x-z,x-y)$，则
  \[
    \nabla\times\vec F=(0,0,2).
  \]
  题设方向从 $z$ 轴正向看为顺时针，对应法向量的 $z$ 分量为负。由 Stokes 公式，
  \[
    I=\iint_D 2(-1)\,dx\,dy=-2\pi,
  \]
  其中 $D:x^2+y^2\le1$。
\end{enumerate}

\paragraph{三、解答题}
\begin{enumerate}
  \item 记 $P=x\varphi(y)$，$Q=x^2y$。路径无关给出
  \[
    P_y=Q_x,\qquad x\varphi'(y)=2xy.
  \]
  故 $\varphi'(y)=2y$，由 $\varphi(0)=1$ 得 $\varphi(y)=y^2+1$。势函数可取
  \[
    \Phi(x,y)=\frac{x^2}{2}(y^2+1),
  \]
  所求积分为 $\Phi(1,1)-\Phi(0,0)=1$。
  \item 令 $t=-x$，则
  \[
    \sum_{n=2}^\infty(-1)^n\frac{x^n}{n(n-1)}
    =\sum_{n=2}^\infty\frac{t^n}{n(n-1)}.
  \]
  当 $|t|<1$ 时，
  \[
    \sum_{n=2}^\infty\frac{t^n}{n(n-1)}
    =(1-t)\ln(1-t)+t.
  \]
  因而
  \[
    S(x)=(1+x)\ln(1+x)-x.
  \]
  端点 $x=1$ 绝对收敛，$x=-1$ 也收敛，故收敛域为 $[-1,1]$。
  \item 齐次方程特征根为 $1,-2$。设特解为 $A\sin x+B\cos x$，代入得
  \[
    A=\frac35,\qquad B=\frac15.
  \]
  通解为
  \[
    y=C_1e^x+C_2e^{-2x}+\frac35\sin x+\frac15\cos x.
  \]
\end{enumerate}

\paragraph{四、应用题}
椭圆可参数化为 $x=2\cos t,\ y=\sin t$。点到直线距离为
\[
  d=\frac{|2x+3y-6|}{\sqrt{13}}.
\]
在椭圆上有 $2x+3y=4\cos t+3\sin t\le5<6$，故距离最小时需使 $4\cos t+3\sin t$ 最大。取
\[
  \cos t=\frac45,\qquad \sin t=\frac35,
\]
得所求点为
\[
  \left(\frac85,\frac35\right).
\]

\paragraph{五、证明题}
设余项
\[
  R_N(x)=\sum_{n=N}^\infty x^2e^{-nx}
  =\frac{x^2e^{-Nx}}{1-e^{-x}},\qquad x>0.
\]
函数 $h(x)=x^2/(e^x-1)$ 在 $(0,+\infty)$ 上有界，且当 $x\to0^+$ 与 $x\to+\infty$ 时均趋于 $0$。
给定 $\varepsilon>0$，先取 $\delta>0$ 使 $0<x<\delta$ 时 $h(x)<\varepsilon$。
再对 $x\ge\delta$ 用
\[
  R_N(x)=h(x)e^{-(N-1)x}\le M e^{-(N-1)\delta},
\]
取 $N$ 足够大即可使该项小于 $\varepsilon$。故 $R_N(x)\to0$ 关于 $x\in(0,+\infty)$ 一致成立，原级数一致收敛。

\section{2022级工科数学分析下 A卷}

\subsection{资料来源}
\begin{itemize}
  \item 2022级A卷无答案.pdf（试题，扫描型 PDF，已按页面人工转写）。
  \item 2022级A卷学生答案.pdf（学生答案，扫描型 PDF；能辨认处用于校对，缺口处按教材方法补写解析）。
\end{itemize}

\subsection{试题}

\paragraph*{一、填空题（共 5 题，每题 2 分，共 10 分）}
\begin{enumerate}
  \item 微分方程 $y''+2y'+6y=0$ 的通解为 \underline{\hspace{4cm}}。
  \item 设函数 $u=2x^2+2y^2+3z^2$，则 $\operatorname{div}(\operatorname{grad}u)=$ \underline{\hspace{3cm}}。
  \item 设 $\Gamma$ 为 $y^2=2x$ 上从原点到 $(2,2)$ 的一段，则第一类曲线积分
  \[
    \int_\Gamma y\,ds=
  \]
  \underline{\hspace{3cm}}。
  \item 参数曲线
  \[
    \begin{cases}
      x=t-\cos t,\\
      y=\sin t,\\
      z=2t
    \end{cases}
  \]
  在 $t=0$ 对应的点处的切线方程为 \underline{\hspace{5cm}}。
  \item 设周期为 $2\pi$ 的函数
  \[
    f(x)=
    \begin{cases}
      -1, & -\pi<x\le 0,\\
      1+x, & 0<x\le \pi,
    \end{cases}
  \]
  则 $f(x)$ 的傅里叶级数在 $x=\pi$ 处收敛于 \underline{\hspace{3cm}}。
\end{enumerate}

\paragraph*{二、选择题（共 5 题，每题 2 分，共 10 分）}
\begin{enumerate}
  \item 关于未知函数 $y$ 的微分方程
  \[
    (y-\cos^2 x)\,dx+e^x\,dy=0
  \]
  是（\quad）。

  \begin{tabular}{llll}
    A. 可分离变量方程； & B. 一阶线性非齐次方程； &
    C. 一阶线性齐次方程； & D. 非线性方程。
  \end{tabular}

  \item 若二元函数 $f(x,y)$ 满足
  \[
    \lim_{(x,y)\to(0,0)}
    \frac{f(x,y)-f(0,0)+3x-5y}{\sqrt{x^2+y^2}}=0,
  \]
  则（\quad）。

  \begin{tabular}{ll}
    A. $df(0,0)=0$； & B. $df(0,0)=3dx-5dy$；\\
    C. $df(0,0)=-3dx+5dy$； & D. $df(0,0)$ 不存在。
  \end{tabular}

  \item 设函数 $f(x,y)$ 与 $\varphi(x,y)$ 均可微，且 $\varphi_y(x,y)\ne0$。若 $(x_0,y_0)$ 是 $f(x,y)$ 在约束条件
  $\varphi(x,y)=0$ 下的一个极值点，则下列选项正确的是（\quad）。

  \begin{tabular}{ll}
    A. 若 $f_x(x_0,y_0)=0$，则 $f_y(x_0,y_0)=0$； &
    B. 若 $f_x(x_0,y_0)=0$，则 $f_y(x_0,y_0)\ne0$；\\
    C. 若 $f_x(x_0,y_0)\ne0$，则 $f_y(x_0,y_0)=0$； &
    D. 若 $f_x(x_0,y_0)\ne0$，则 $f_y(x_0,y_0)\ne0$。
  \end{tabular}

  \item 函数 $\ln(1+x)$ 在 $x=0$ 处的泰勒展开式正确的是（\quad）。

  \begin{tabular}{ll}
    A. $\displaystyle \ln(1+x)=\sum_{n=1}^{\infty}\frac{(-1)^n}{n}x^n,\ x\in(-1,1]$； &
    B. $\displaystyle \ln(1+x)=\sum_{n=1}^{\infty}\frac{(-1)^{n-1}}{n}x^n,\ x\in(-1,1]$；\\
    C. $\displaystyle \ln(1+x)=-\sum_{n=1}^{\infty}\frac{1}{n}x^n,\ x\in[-1,1)$； &
    D. $\displaystyle \ln(1+x)=\sum_{n=1}^{\infty}\frac{1}{n}x^n,\ x\in[-1,1)$。
  \end{tabular}

  \item 使得级数
  \[
    \sum_{n=1}^{\infty}\frac{(-1)^n}{n^p}
  \]
  条件收敛的常数 $p$ 的取值范围是（\quad）。

  \begin{tabular}{llll}
    A. $p\le0$； & B. $0<p\le1$； & C. $0<p<1$； & D. $p>1$。
  \end{tabular}
\end{enumerate}

\paragraph*{三、计算题（共 3 题，每题 10 分，共 30 分）}
\begin{enumerate}
  \item 设
  \[
    z=\frac{1}{x}f(x+y,x-y)+g(xy),
  \]
  其中函数 $f$ 有二阶连续的偏导数，且 $g$ 二阶可导，计算 $\dfrac{\partial z}{\partial x}$ 和
  $\dfrac{\partial^2z}{\partial x\partial y}$。

  \item 计算累次积分
  \[
    \int_{-\sqrt2}^{0}dx\int_{-x}^{\sqrt{4-x^2}}\sqrt{x^2+y^2}\,dy
    +\int_{0}^{2}dx\int_{\sqrt{2x-x^2}}^{\sqrt{4-x^2}}\sqrt{x^2+y^2}\,dy.
  \]

  \item 计算锥面 $z=\sqrt{x^2+y^2}$ 被柱面 $x^2+y^2=4x$ 截下的部分锥面面积。
\end{enumerate}

\paragraph*{四、综合题（共 3 题，每题 10 分，共 30 分）}
\begin{enumerate}
  \item 设 $\Sigma$ 是
  \[
    \frac{x^2}{a^2}+\frac{y^2}{b^2}+\frac{z^2}{c^2}=1\quad (z\ge0)
  \]
  的上侧，且 $a,b,c$ 都是正实数。计算第二类曲面积分
  \[
    \iint_\Sigma (x^2-y\sin x)\,dy\,dz+(y^2-z^2)\,dz\,dx+(z^2-x^2)\,dx\,dy.
  \]

  \item 计算曲线积分
  \[
    \int_\Gamma
    \frac{(x+y)\,dx-(x-y)\,dy}{x^2+y^2},
  \]
  其中曲线 $\Gamma$ 是从点 $A(-1,0)$ 到点 $B(1,0)$ 的一条不经过原点的光滑曲线 $y=f(x)$，
  $-1\le x\le1$。

  \item 求幂级数
  \[
    \sum_{n=0}^{\infty}(n+1)x^{2n+1}
  \]
  的收敛域及和函数。
\end{enumerate}

\paragraph*{五、证明题（共 1 题，每题 10 分，共 10 分）}
证明函数项级数
\[
  \sum_{n=1}^{\infty}\sin\frac{n!x^n}{x^2+n^n}
\]
在区间 $[-2,2]$ 上一致收敛。

\paragraph*{六、应用题（共 1 题，每题 10 分，共 10 分）}
制作一个体积为 $2$ 的长方体盒子，利用拉格朗日乘数法求出长、宽和高分别为多少时才能使其表面积最小。

\subsection{答案与解析}

\paragraph*{一、填空题}
\begin{enumerate}
  \item 特征方程为 $r^2+2r+6=0$，根为 $r=-1\pm\sqrt5\,i$，故
  \[
    y=e^{-x}\left(C_1\cos\sqrt5\,x+C_2\sin\sqrt5\,x\right).
  \]
  \item $\operatorname{div}(\operatorname{grad}u)=\Delta u=u_{xx}+u_{yy}+u_{zz}=4+4+6=14$。
  \item 令 $y=t,\ x=\dfrac{t^2}{2}$，$0\le t\le2$，则 $ds=\sqrt{1+t^2}\,dt$，所以
  \[
    \int_\Gamma y\,ds=\int_0^2 t\sqrt{1+t^2}\,dt
    =\frac{(1+t^2)^{3/2}}{3}\bigg|_0^2
    =\frac{5\sqrt5-1}{3}.
  \]
  \item 当 $t=0$ 时点为 $(-1,0,0)$，切向量为 $(1,1,2)$，故切线方程为
  \[
    \frac{x+1}{1}=\frac{y}{1}=\frac{z}{2}.
  \]
  \item 傅里叶级数在跳跃点收敛于左右极限的平均值。因
  \[
    f(\pi-0)=1+\pi,\qquad f(\pi+0)=f(-\pi+0)=-1,
  \]
  所以收敛值为 $\dfrac{\pi}{2}$。
\end{enumerate}

\paragraph*{二、选择题}
\begin{enumerate}
  \item 方程可化为 \(y'+e^{-x}y=e^{-x}\cos^2x\)，是一阶线性非齐次方程。选 B。
  \item 由可微定义，线性主部为 \(-3x+5y\)，故 \(df(0,0)=-3\,dx+5\,dy\)。选 C。
  \item 约束极值必要条件给出 \(f_x\varphi_y-f_y\varphi_x=0\)。因 \(\varphi_y\ne0\)，若 \(f_x\ne0\) 则必有 \(f_y\ne0\)。选 D。
  \item \(\displaystyle \ln(1+x)=\sum_{n=1}^{\infty}\frac{(-1)^{n-1}}{n}x^n\)，收敛区间为 $(-1,1]$。选 B。
  \item 交错 $p$ 级数在 $p>0$ 时收敛，在 $p>1$ 时绝对收敛，所以条件收敛范围为 $0<p\le1$。选 B。
\end{enumerate}

\paragraph*{三、计算题}
\begin{enumerate}
  \item 记 $u=x+y,\ v=x-y$，并把 $f_1,f_2,f_{11},f_{22}$ 均理解为在 $(u,v)$ 处的取值，把 $g',g''$
  理解为在 $xy$ 处的取值。先对 $x$ 求导：
  \[
    \frac{\partial z}{\partial x}
    =-\frac{f}{x^2}+\frac{f_1+f_2}{x}+y\,g'(xy).
  \]
  再对 $y$ 求导，利用 $u_y=1,\ v_y=-1$，得
  \[
    \frac{\partial^2z}{\partial x\partial y}
    =-\frac{f_1-f_2}{x^2}+\frac{f_{11}-f_{22}}{x}+g'(xy)+xy\,g''(xy).
  \]

  \item 第一部分区域为极坐标下
  \[
    \frac{\pi}{2}\le\theta\le\frac{3\pi}{4},\qquad 0\le r\le2,
  \]
  第二部分区域为
  \[
    0\le\theta\le\frac{\pi}{2},\qquad 2\cos\theta\le r\le2.
  \]
  因 $\sqrt{x^2+y^2}=r$，面积元为 $r\,dr\,d\theta$，所以原积分为
  \[
    \int_0^{\pi/2}\int_{2\cos\theta}^{2}r^2\,dr\,d\theta
    +\int_{\pi/2}^{3\pi/4}\int_0^2r^2\,dr\,d\theta
    =2\pi-\frac{16}{9}.
  \]

  \item 锥面写成 $z=\sqrt{x^2+y^2}$ 时，
  \[
    dS=\sqrt{1+z_x^2+z_y^2}\,dx\,dy=\sqrt2\,dx\,dy.
  \]
  投影区域为圆 $x^2+y^2\le4x$，即半径为 $2$ 的圆盘，面积为 $4\pi$。故锥面面积为
  \[
    S=\sqrt2\cdot4\pi=4\sqrt2\,\pi.
  \]
\end{enumerate}

\paragraph*{四、综合题}
\begin{enumerate}
  \item 记
  \[
    P=x^2-y\sin x,\qquad Q=y^2-z^2,\qquad R=z^2-x^2.
  \]
  用底面
  \[
    \Sigma_0:\ z=0,\quad \frac{x^2}{a^2}+\frac{y^2}{b^2}\le1
  \]
  补成上半椭球体 $\Omega$ 的外侧闭曲面，底面取向向下。由高斯公式，
  \[
    I_\Sigma+I_{\Sigma_0}
    =\iiint_\Omega (P_x+Q_y+R_z)\,dV
    =\iiint_\Omega (2x-y\cos x+2y+2z)\,dV.
  \]
  上半椭球关于 $x,y$ 对称，含 $x,y$ 的奇函数项积分为 $0$，故
  \[
    I_\Sigma+I_{\Sigma_0}=\iiint_\Omega 2z\,dV
    =\frac{\pi ab c^2}{2}.
  \]
  底面向下时
  \[
    I_{\Sigma_0}=\iint_{\frac{x^2}{a^2}+\frac{y^2}{b^2}\le1}x^2\,dx\,dy
    =\frac{\pi a^3b}{4}.
  \]
  因而
  \[
    I_\Sigma=\frac{\pi ab c^2}{2}-\frac{\pi a^3b}{4}.
  \]

  \item 积分可拆成
  \[
    \frac{(x+y)\,dx-(x-y)\,dy}{x^2+y^2}
    =\frac{x\,dx+y\,dy}{x^2+y^2}
    +\frac{y\,dx-x\,dy}{x^2+y^2}
    =d\ln r-d\theta.
  \]
  两端点到原点的距离都为 $1$，故 $d\ln r$ 的积分为 $0$。剩下的值由曲线绕过原点的方向决定。若
  $f(0)>0$，曲线从上半平面绕过原点，$\theta$ 可由 $\pi$ 连续变到 $0$，积分为 $\pi$；若
  $f(0)<0$，曲线从下半平面绕过原点，积分为 $-\pi$。因此
  \[
    \int_\Gamma
    \frac{(x+y)\,dx-(x-y)\,dy}{x^2+y^2}
    =
    \begin{cases}
      \pi, & f(0)>0,\\
      -\pi, & f(0)<0.
    \end{cases}
  \]

  \item 令 $t=x^2$，则
  \[
    \sum_{n=0}^{\infty}(n+1)x^{2n+1}
    =x\sum_{n=0}^{\infty}(n+1)t^n.
  \]
  当 $|t|<1$ 即 $|x|<1$ 时，
  \[
    \sum_{n=0}^{\infty}(n+1)t^n=\frac{1}{(1-t)^2},
  \]
  所以和函数为
  \[
    S(x)=\frac{x}{(1-x^2)^2},\qquad |x|<1.
  \]
  当 $x=\pm1$ 时通项不趋于 $0$，故发散。因此收敛域为 $(-1,1)$。
\end{enumerate}

\paragraph*{五、证明题}
对任意 $x\in[-2,2]$，有
\[
  \left|\sin\frac{n!x^n}{x^2+n^n}\right|
  \le
  \left|\frac{n!x^n}{x^2+n^n}\right|
  \le
  \frac{n!\,2^n}{n^n}.
\]
令 $M_n=\dfrac{n!\,2^n}{n^n}$，则
\[
  \frac{M_{n+1}}{M_n}
  =2\left(\frac{n}{n+1}\right)^n\longrightarrow \frac{2}{e}<1.
\]
故 $\sum M_n$ 收敛。由魏尔斯特拉斯判别法，原函数项级数在 $[-2,2]$ 上一致收敛。

\paragraph*{六、应用题}
设长方体长、宽、高分别为 $x,y,z>0$，体积约束为 $xyz=2$，闭盒子的表面积为
\[
  S=2(xy+xz+yz).
\]
设
\[
  L=2(xy+xz+yz)+\lambda(xyz-2).
\]
由拉格朗日乘数法，
\[
  2(y+z)+\lambda yz=0,\qquad
  2(x+z)+\lambda xz=0,\qquad
  2(x+y)+\lambda xy=0.
\]
两两相减得 $x=y=z$。再由 $xyz=2$，得
\[
  x=y=z=\sqrt[3]{2}.
\]
因此长、宽和高都等于 $\sqrt[3]{2}$ 时表面积最小。

\section{2022级工科数学分析下 B卷}

\subsection{试题}

\paragraph*{一、填空题：共 5 题，每题 2 分，共 10 分。}
\begin{enumerate}
  \item 微分方程 $y''+3y'+2y=0$ 的通解为\underline{\hspace{4cm}}。
  \item 函数 $u=2xy-z^2$ 在点 $(2,-1,-1)$ 处方向导数的最大值为\underline{\hspace{4cm}}。
  \item 设 $\Gamma$ 是圆周 $x^2+y^2=1$ 在第一象限的部分，则第一类曲线积分
  \[
    \int_\Gamma (x+y)^2\,ds=\underline{\hspace{4cm}}.
  \]
  \item 级数 $\displaystyle \sum_{n=0}^{\infty}(-1)^n\dfrac{2n+3}{(2n+1)!}$ 的和为\underline{\hspace{4cm}}。
  \item 设周期为 $2\pi$ 的函数
  \[
    f(x)=
    \begin{cases}
      x^2, & -\pi<x\le 0,\\
      0, & 0<x\le \pi,
    \end{cases}
  \]
  则 $f(x)$ 的傅里叶（Fourier）级数在 $x=-\pi$ 处收敛于\underline{\hspace{4cm}}。
\end{enumerate}

\paragraph*{二、选择题：共 5 题，每题 2 分，共 10 分。}
\begin{enumerate}
  \item 下列微分方程中，属于二阶线性常微分方程的是（\quad）
  \begin{enumerate}
    \item $y''+x^3y'+(\ln x)y=\tan x$；
    \item $(x^2+y^2)\,dy+y^2\,dx=0$；
    \item $(y'')^2+x^2y=1$；
    \item $y''+2\ln y=2x$。
  \end{enumerate}
  \item 已知函数 $f(x)$ 具有二阶连续导函数，且 $f(x)>0,\ f'(0)=0$（$f$ 在 $0$ 处的导数为 $0$）。则函数 $z=f(x)\ln f(y)$ 在点 $(0,0)$ 处取得极小值的一个充分条件是（\quad）
  \begin{enumerate}
    \item $f(0)>1,\ f''(0)>0$；
    \item $f(0)>1,\ f''(0)<0$；
    \item $f(0)<1,\ f''(0)>0$；
    \item $f(0)<1,\ f''(0)<0$。
  \end{enumerate}
  \item 函数 $f(x,y)=1+x+y$ 在区域 $\{(x,y)\mid x^2+y^2\le 1\}$ 上的最大值与最小值之积是（\quad）
  \begin{enumerate}
    \item $-1$；
    \item $1$；
    \item $3-\sqrt2$；
    \item $1+\sqrt2$。
  \end{enumerate}
  \item 设 $\Gamma$ 是以 $(1,1),(-1,1),(-1,-1),(1,-1)$ 为顶点的正方形边界，则第一类曲线积分
  \[
    \oint_\Gamma \frac{x+y+1}{|x|+|y|}\,ds
  \]
  等于（\quad）
  \begin{enumerate}
    \item $0$；
    \item $8$；
    \item $4\ln 2$；
    \item $8\ln 2$。
  \end{enumerate}
  \item 设 $a$ 为常数，则级数
  \[
    \sum_{n=1}^{\infty}\left(\frac{\sin(an)}{n^2}-\frac1{\sqrt n}\right)
  \]
  （\quad）
  \begin{enumerate}
    \item 发散；
    \item 绝对收敛；
    \item 条件收敛；
    \item 敛散性与 $a$ 的取值相关。
  \end{enumerate}
\end{enumerate}

\paragraph*{三、计算题：共 3 题，每题 10 分，共 30 分。}
\begin{enumerate}
  \item 设
  \[
    g(x,y)=f\left(xy,\frac12(x^2-y^2)\right),
  \]
  其中 $f(u,v)$ 具有连续的二阶偏导数，且满足
  \[
    \frac{\partial^2 f}{\partial u^2}+\frac{\partial^2 f}{\partial v^2}=1.
  \]
  计算二阶偏导数 $\displaystyle \frac{\partial^2 g}{\partial x^2}+\frac{\partial^2 g}{\partial y^2}$。
  \item 计算累次积分
  \[
    \int_0^1 dx\int_0^x dy\int_0^y \frac{\sin z}{(1-z)^2}\,dz.
  \]
  \item 计算二重积分
  \[
    \iint_D \sqrt{x^2+y^2}\,dx\,dy,
  \]
  其中 $D$ 是由 $x=\sqrt{2y-y^2}$ 与 $y=x$ 所围成的闭区域。
\end{enumerate}

\paragraph*{四、综合题：共 3 题，每题 10 分，共 30 分。}
\begin{enumerate}
  \item 设 $\Sigma$ 是曲面 $z=1-x^2-y^2\ (z\ge 0)$ 的上侧，计算第二型曲面积分
  \[
    \iint_\Sigma 2x^3\,dy\,dz+2y^3\,dz\,dx+3(z^2-1)\,dx\,dy.
  \]
  \item 计算曲线积分
  \[
    \oint_L yz\,dx+3zx\,dy-xy\,dz,
  \]
  其中 $L$ 是曲线
  \[
    \begin{cases}
      x^2+y^2=2y,\\
      y-z=1,
    \end{cases}
  \]
  其方向从 $z$ 轴正向往 $z$ 轴负向看去为逆时针方向。
  \item 将函数
  \[
    f(x)=\frac{x}{x^2-5x+6}
  \]
  展开成 $x-1$ 的幂级数，并指出其收敛域。
\end{enumerate}

\paragraph*{五、证明题：共 1 题，每题 10 分，共 10 分。}
证明函数项级数
\[
  \sum_{n=1}^{\infty}\arctan\frac{2x}{x^2+n^3}
\]
在区间 $(-\infty,+\infty)$ 上一致收敛。

\paragraph*{六、解答题：共 1 题，每题 10 分，共 10 分。}
限制点 $(x,y)$ 在圆周 $(x-1)^2+y^2=1$ 上变化时，求函数 $f(x,y)=xy$ 的最小值与最大值。

\subsection{答案与解析}

\paragraph*{一、填空题}
\begin{enumerate}
  \item 特征方程 $r^2+3r+2=0$，得 $r=-1,-2$，故通解为
  \[
    y=C_1e^{-x}+C_2e^{-2x}.
  \]
  \item $\nabla u=(2y,2x,-2z)$，在 $(2,-1,-1)$ 处为 $(-2,4,2)$，方向导数最大值为
  \[
    |\nabla u|=\sqrt{(-2)^2+4^2+2^2}=2\sqrt6.
  \]
  \item 令 $x=\cos t,\ y=\sin t,\ 0\le t\le \dfrac{\pi}{2}$，则
  \[
    \int_\Gamma (x+y)^2\,ds
    =\int_0^{\pi/2}(\cos t+\sin t)^2\,dt
    =\frac{\pi}{2}+1.
  \]
  \item
  \[
    \sum_{n=0}^{\infty}(-1)^n\frac{2n+3}{(2n+1)!}
    =
    \sum_{n=0}^{\infty}\frac{(-1)^n}{(2n)!}
    +2\sum_{n=0}^{\infty}\frac{(-1)^n}{(2n+1)!}
    =\cos 1+2\sin 1.
  \]
  \item 在 $x=-\pi$ 处，周期延拓函数的右极限为 $\pi^2$，左极限为 $0$，故 Fourier 级数收敛于
  \[
    \frac{\pi^2+0}{2}=\frac{\pi^2}{2}.
  \]
\end{enumerate}

\paragraph*{二、选择题}
\begin{enumerate}
  \item 选 A。A 为二阶线性非齐次常微分方程；B 为一阶方程；C 含 $(y'')^2$，D 含 $\ln y$，均非二阶线性方程。
  \item 选 A。由 $f'(0)=0$，点 $(0,0)$ 为驻点。Hessian 矩阵对角元为 $f''(0)\ln f(0)$ 与 $f''(0)$，交叉项为 $0$；当 $f(0)>1,\ f''(0)>0$ 时正定，故取得极小值。
  \item 选 A。在线性函数 $1+x+y$ 在单位圆盘上，最大值为 $1+\sqrt2$，最小值为 $1-\sqrt2$，二者乘积为 $-1$。
  \item 选 D。沿正方形四条边分段计算，位于上边与右边的积分各为 $4\ln2$，下边与左边的积分和为 $0$，故总积分为 $8\ln2$。
  \item 选 A。$\sum \sin(an)/n^2$ 绝对收敛，而 $\sum 1/\sqrt n$ 发散，故原级数发散。
\end{enumerate}

\paragraph*{三、计算题}
\begin{enumerate}
  \item 记 $u=xy,\ v=\dfrac12(x^2-y^2)$，并用 $f_1,f_2$ 表示 $f$ 对第一、第二个变量的偏导数。则
  \[
    g_x=yf_1+xf_2,\qquad g_y=xf_1-yf_2.
  \]
  继续求导得
  \[
    g_{xx}=y^2f_{11}+2xyf_{12}+x^2f_{22}+f_2,
  \]
  \[
    g_{yy}=x^2f_{11}-2xyf_{12}+y^2f_{22}-f_2.
  \]
  因此
  \[
    g_{xx}+g_{yy}=(x^2+y^2)(f_{11}+f_{22})=x^2+y^2.
  \]
  \item 积分区域为 $0\le z\le y\le x\le 1$。交换积分次序得
  \[
    \int_0^1\frac{\sin z}{(1-z)^2}
    \left(\int_z^1dy\int_y^1dx\right)dz
    =
    \int_0^1\frac{\sin z}{(1-z)^2}\cdot\frac{(1-z)^2}{2}\,dz.
  \]
  故原式为
  \[
    \frac12\int_0^1\sin z\,dz=\frac{1-\cos 1}{2}.
  \]
  \item 曲线 $x=\sqrt{2y-y^2}$ 等价于 $x^2+(y-1)^2=1$，在极坐标下边界为 $r=2\sin\theta$；直线 $y=x$ 为 $\theta=\dfrac{\pi}{4}$。故
  \[
    D:\quad 0\le \theta\le \frac{\pi}{4},\quad 0\le r\le 2\sin\theta.
  \]
  于是
  \[
    \iint_D\sqrt{x^2+y^2}\,dx\,dy
    =
    \int_0^{\pi/4}\int_0^{2\sin\theta}r^2\,dr\,d\theta
    =
    \frac{16-10\sqrt2}{9}.
  \]
\end{enumerate}

\paragraph*{四、综合题}
\begin{enumerate}
  \item 用平面 $\Sigma_1:z=0,\ x^2+y^2\le 1$ 补成闭曲面，取外侧方向。由高斯公式，
  \[
    \iint_{\Sigma+\Sigma_1}2x^3\,dy\,dz+2y^3\,dz\,dx+3(z^2-1)\,dx\,dy
    =
    \iiint_\Omega (6x^2+6y^2+6z)\,dV=2\pi.
  \]
  在补面 $\Sigma_1$ 上取下侧，积分为
  \[
    I_{\Sigma_1}=-\iint_{x^2+y^2\le 1}3(z^2-1)\,dx\,dy=3\pi.
  \]
  因此所求曲面积分为
  \[
    2\pi-3\pi=-\pi.
  \]
  \item 按题设方向可取参数
  \[
    x=\cos t,\qquad y=1+\sin t,\qquad z=\sin t,\qquad 0\le t\le 2\pi.
  \]
  代入得
  \[
    \oint_L yz\,dx+3zx\,dy-xy\,dz
    =
    \int_0^{2\pi}\left[-(1+\sin t)\sin^2t+3\sin t\cos^2t-(1+\sin t)\cos^2t\right]dt.
  \]
  利用一个周期内奇函数项积分为 $0$，得
  \[
    -\int_0^{2\pi}\sin^2t\,dt-\int_0^{2\pi}\cos^2t\,dt=-2\pi.
  \]
  \item 分解为
  \[
    \frac{x}{x^2-5x+6}=\frac{3}{x-3}-\frac{2}{x-2}.
  \]
  令 $t=x-1$，则
  \[
    \frac{3}{x-3}=-\frac32\cdot\frac1{1-t/2}
    =-\frac32\sum_{n=0}^{\infty}\left(\frac{t}{2}\right)^n,\quad |t|<2,
  \]
  \[
    -\frac{2}{x-2}=\frac2{1-t}
    =2\sum_{n=0}^{\infty}t^n,\quad |t|<1.
  \]
  因此
  \[
    f(x)=\sum_{n=0}^{\infty}\left(2-\frac{3}{2^{n+1}}\right)(x-1)^n,
    \qquad |x-1|<1.
  \]
  收敛域为 $(0,2)$。
\end{enumerate}

\paragraph*{五、证明题}
对任意 $x\in\mathbb R$，由 $|\arctan u|\le |u|$ 以及
\[
  x^2+n^3\ge 2|x|n^{3/2}
\]
可得
\[
  \left|\arctan\frac{2x}{x^2+n^3}\right|
  \le
  \frac{2|x|}{x^2+n^3}
  \le
  \frac1{n^{3/2}}.
\]
而级数 $\sum_{n=1}^{\infty}n^{-3/2}$ 收敛，故由 Weierstrass 判别法知
\[
  \sum_{n=1}^{\infty}\arctan\frac{2x}{x^2+n^3}
\]
在 $(-\infty,+\infty)$ 上一致收敛。

\paragraph*{六、解答题}
令约束函数
\[
  \varphi(x,y)=(x-1)^2+y^2-1=0.
\]
对 $F(x,y,\lambda)=xy+\lambda\varphi(x,y)$，有
\[
  y+2\lambda(x-1)=0,\qquad x+2\lambda y=0.
\]
由约束得 $x^2+y^2=2x$。当 $\lambda=0$ 时得点 $(0,0)$，函数值为 $0$，不是最值点。设 $y\ne0$，由第二个方程得
\[
  \lambda=-\frac{x}{2y}.
\]
代入第一个方程，得
\[
  y^2=x(x-1).
\]
再与 $y^2=2x-x^2$ 联立，得
\[
  x(x-1)=2x-x^2,\qquad x=\frac32,\qquad y=\pm\frac{\sqrt3}{2}.
\]
因此
\[
  f\left(\frac32,\frac{\sqrt3}{2}\right)=\frac{3\sqrt3}{4}
  \]
为最大值，
\[
  f\left(\frac32,-\frac{\sqrt3}{2}\right)=-\frac{3\sqrt3}{4}
  \]
为最小值。

\section{2023级工科数学分析（二）期末考试 A 卷}

\subsection{试题}

\paragraph{一、填空题（共 5 小题，每小题 2 分，共 10 分）}
\begin{enumerate}
  \item 方程 $y'-\dfrac{y}{x}=x^2$ 的通解为 \underline{\hspace{4cm}}。
  \item 设 $z=\ln(x^2+y^2)$，$\vec n=\left(\cos\dfrac{\pi}{3},\sin\dfrac{\pi}{3}\right)$。则
  \[
    \left.\frac{\partial z}{\partial \vec n}\right|_{(1,1)}=\underline{\hspace{4cm}}.
  \]
  \item 设 $\Omega=\{(x,y,z):x^2+y^2+z^2\le 1\}$，则
  \[
    \iiint_\Omega (x^2+y^2+z^2)\,dV=\underline{\hspace{4cm}}.
  \]
  \item 设 $C$ 是从点 $A=(3,2,1)$ 到点 $B=(0,0,0)$ 的直线段，则
  \[
    \int_C x^3\,dx+y^3\,dy+z^3\,dz=\underline{\hspace{4cm}}.
  \]
  \item 设 $f(x)$ 是以 $2\pi$ 为周期的函数，它在 $[-\pi,\pi)$ 上的表达式是
  \[
    f(x)=
    \begin{cases}
      x+1, & -\pi\le x<0,\\
      x^2, & 0\le x<\pi.
    \end{cases}
  \]
  若它的 Fourier 级数的和函数是 $S(x)$，则 $S(-\pi)=\underline{\hspace{4cm}}$。
\end{enumerate}

\paragraph{二、计算题（共 4 小题，每小题 9 分，共 36 分）}
\begin{enumerate}
  \item 求方程 $y''-y=\sin^2x$ 的通解。
  \item 设 $z=\arctan(x+y)$，$x=2u-v^2$，$y=u^2v$，求 $\dfrac{\partial z}{\partial u}$，$\dfrac{\partial z}{\partial v}$。
  \item 设 $D$ 是以 $(0,0)$，$(\pi,-\pi)$，$(\pi,\pi)$ 为顶点的三角形区域，求
  \[
    \iint_D \cos(x-y)\sin(x+y)\,dx\,dy.
  \]
  \item 设 $\Omega\subset\mathbb R^3$ 是由曲面 $z=5-x^2-y^2$ 和 $z+1=\sqrt{x^2+y^2}$ 所围成的空间立体，且 $\Omega$ 内布满某种流体，其流速为
  \[
    \vec F=\frac{\vec r}{\|\vec r\|^3},\qquad \vec r=(x,y,z).
  \]
  求单位时间内流体通过 $\Omega$ 的表面 $\partial\Omega$ 的流量。
\end{enumerate}

\paragraph{三、解答题（共 3 小题，每小题 9 分，共 27 分）}
\begin{enumerate}
  \item 求级数 $\displaystyle \sum_{n=1}^\infty \frac{2n-1}{2^n}$ 的和。
  \item 求曲线积分
  \[
    \oint_L y\,dx+z\,dy+x\,dz,
  \]
  其中 $L$ 为圆周
  \[
    \begin{cases}
      x^2+y^2+z^2=a^2,\quad a>0,\\
      x+y+z=0,
    \end{cases}
  \]
  且从 $x$ 轴正向看去，这个圆周取逆时针方向。
  \item 给定
  \[
    du=\frac{(x+2y)\,dx+y\,dy}{(x+y)^2},
  \]
  试验证势函数 $u$ 的存在性，并求出 $u$。
\end{enumerate}

\paragraph{四、应用题（9 分）}
求曲线
\[
  \begin{cases}
    z=x^2+y^2,\\
    y=\dfrac1x
  \end{cases}
\]
上到 $Oxy$ 平面距离最近的点。

\paragraph{五、证明题（共 2 小题，每小题 9 分，共 18 分）}
\begin{enumerate}
  \item 证明曲线
  \[
    \begin{cases}
      x^2+y^2+z^2=6,\\
      x+y+z=0
    \end{cases}
  \]
  在点 $(1,-2,1)$ 处的法平面方程为 $x-z=0$。
  \item 证明级数
  \[
    \sum_{n=3}^\infty \ln\left(1+\frac{x}{n\ln^2 n}\right)
  \]
  在 $[-2,2]$ 内一致收敛。
\end{enumerate}

\subsection{答案与解析}

\paragraph{一、填空题}
\begin{enumerate}
  \item $y=Cx+\dfrac{x^3}{2}$。
  \item $\dfrac{1+\sqrt3}{2}$。
  \item $\dfrac{4\pi}{5}$。
  \item $-\dfrac{49}{2}$。
  \item $\dfrac{\pi^2-\pi+1}{2}$。
\end{enumerate}

\paragraph{二、计算题}
\begin{enumerate}
  \item 因 $\sin^2x=\dfrac{1-\cos2x}{2}$，通解为
  \[
    y=C_1e^x+C_2e^{-x}-\frac12+\frac1{10}\cos2x.
  \]
  \item 令 $s=x+y=2u-v^2+u^2v$，则
  \[
    \frac{\partial z}{\partial u}
    =\frac{2+2uv}{1+s^2},\qquad
    \frac{\partial z}{\partial v}
    =\frac{-2v+u^2}{1+s^2}.
  \]
  \item 作变换 $u=x-y,\ v=x+y$，则 $dx\,dy=\frac12\,du\,dv$，区域变为
  \[
    u\ge0,\quad v\ge0,\quad u+v\le2\pi.
  \]
  因此
  \[
    I=\frac12\int_0^{2\pi}\int_0^{2\pi-u}\cos u\sin v\,dv\,du
    =-\frac{\pi}{2}.
  \]
  \item 该区域包含原点，而 $\vec F=\vec r/\|\vec r\|^3$ 在原点有孤立奇点。去掉原点附近小球并用高斯公式，外边界通量等于小球面通量
  \[
    \iint_{\partial B_\varepsilon}\frac{\vec r}{\|\vec r\|^3}\cdot\vec n\,dS=4\pi.
  \]
  故通过 $\partial\Omega$ 的流量为 $4\pi$。
\end{enumerate}

\paragraph{三、解答题}
\begin{enumerate}
  \item
  \[
    \sum_{n=1}^\infty\frac{2n-1}{2^n}
    =2\sum_{n=1}^\infty\frac{n}{2^n}-\sum_{n=1}^\infty\frac1{2^n}
    =4-1=3.
  \]
  \item 令 $\vec F=(y,z,x)$，则 $\nabla\times\vec F=(-1,-1,-1)$。按题设方向取单位法向
  \[
    \vec n=\frac{(1,1,1)}{\sqrt3}.
  \]
  圆盘半径为 $a$，由 Stokes 公式得
  \[
    \oint_L y\,dx+z\,dy+x\,dz
    =(-1,-1,-1)\cdot\frac{(1,1,1)}{\sqrt3}\,\pi a^2
    =-\sqrt3\,\pi a^2.
  \]
  \item 记
  \[
    P=\frac{x+2y}{(x+y)^2},\qquad Q=\frac{y}{(x+y)^2}.
  \]
  有 $P_y=Q_x=-2y/(x+y)^3$，故势函数存在。对 $P$ 关于 $x$ 积分，得
  \[
    u=\ln|x+y|-\frac{y}{x+y}+C.
  \]
\end{enumerate}

\paragraph{四、应用题}
到 $Oxy$ 平面的距离为 $|z|$。在曲线上 $z=x^2+y^2=x^2+x^{-2}$，故
\[
  z\ge2,
\]
等号在 $x=\pm1$ 时成立。最近点为
\[
  (1,1,2),\qquad (-1,-1,2).
\]

\paragraph{五、证明题}
\begin{enumerate}
  \item 设
  \[
    F=x^2+y^2+z^2-6,\qquad G=x+y+z.
  \]
  在 $(1,-2,1)$ 处，
  \[
    \nabla F=(2,-4,2),\qquad \nabla G=(1,1,1),
  \]
  切向量可取 $\nabla F\times\nabla G=(-6,0,6)$，即平行于 $(1,0,-1)$。法平面垂直于切向量，故
  \[
    (1,0,-1)\cdot(x-1,y+2,z-1)=0,
  \]
  即 $x-z=0$。
  \item 当 $|x|\le2$ 时，各项均有定义。对充分大的 $n$，
  \[
    \left|\ln\left(1+\frac{x}{n\ln^2n}\right)\right|
    \le C\frac{|x|}{n\ln^2n}
    \le \frac{2C}{n\ln^2n}.
  \]
  因 $\sum_{n=3}^\infty 1/(n\ln^2n)$ 收敛，有限个初始项不影响一致收敛，由 Weierstrass 判别法知原级数在 $[-2,2]$ 上一致收敛。
\end{enumerate}

\section{2023级工科数学分析（二）期末考试 B 卷}

\subsection{试题}

\paragraph{一、填空题（共 5 小题，每小题 2 分，共 10 分）}
\begin{enumerate}
  \item 方程 $y'+y\tan x=\cos x$ 的通解为 \underline{\hspace{4cm}}。
  \item 设 $z=e^{xy}$，$\vec n=\left(\cos\dfrac{\pi}{4},\sin\dfrac{\pi}{4}\right)$。则
  \[
    \left.\frac{\partial z}{\partial \vec n}\right|_{(1,1)}=\underline{\hspace{4cm}}.
  \]
  \item
  \[
    \int_0^1 dx\int_0^{\sqrt{1-x^2}}x^3\,dy=\underline{\hspace{4cm}}.
  \]
  \item 设曲线的参数方程为
  \[
    x=at,\qquad y=\frac a2t^2,\qquad z=\frac a3t^3,\qquad 0\le t\le 1,
  \]
  且其密度函数为 $\sqrt{2y/a}$，则曲线的质量为 \underline{\hspace{4cm}}。
  \item 设 $f(x)$ 是以 $2\pi$ 为周期的函数，它在 $[-\pi,\pi)$ 上的表达式是 $f(x)=x$。则 $f(x)$ 的 Fourier 级数中 $\sin 3x$ 的系数是 \underline{\hspace{4cm}}。
\end{enumerate}

\paragraph{二、计算题（共 4 小题，每小题 9 分，共 36 分）}
\begin{enumerate}
  \item 求方程 $y''+y=3xe^x$ 的通解。
  \item 设 $z=\sin\left(xy+\dfrac{x}{y}\right)$，求 $\dfrac{\partial^2z}{\partial x\partial y}$。
  \item 设 $D$ 是由 $x^2+y^2=9$，$x^2+y^2=16$ 以及 $y^2-x^2=1$，$y^2-x^2=9$ 围成的位于第一象限的区域，求
  \[
    \iint_D xy\,dx\,dy.
  \]
  \item 设 $\Omega\subset\mathbb R^3$ 以原点为内点，且边界曲面 $\partial\Omega$ 光滑，又设边界曲面单位外法向为 $\vec n$，向量场
  \[
    \vec F=\frac{\vec r}{\|\vec r\|^3},\qquad \vec r=(x,y,z).
  \]
  求
  \[
    \iint_{\partial\Omega}\vec F\cdot\vec n\,ds.
  \]
\end{enumerate}

\paragraph{三、解答题（共 3 小题，每小题 9 分，共 27 分）}
\begin{enumerate}
  \item 求级数 $\displaystyle \sum_{n=1}^\infty \frac1{4n(4n+2)}$ 的和。
  \item 如图，设有向曲面
  \[
    \Sigma=\{(x,y,z):x+y+z=1,\ x\ge 0,\ y\ge 0,\ z\ge 0\},
  \]
  且其方向与 $(1,1,1)$ 同向。求力场 $\vec F=(y^2,z^2,x^2)$ 绕 $\Sigma$ 正向边界 $\partial\Sigma$ 一周所做的功。
  \item 已知 $\phi(\pi)=1$，试确定 $\phi(x)$，使得曲线积分
  \[
    I=\int_A^B \left[\sin x-\phi(x)\right]\frac{y}{x}\,dx+\phi(x)\,dy
  \]
  与路径无关，并求当 $A,B$ 两点分别为 $(1,0)$，$(\pi,\pi)$ 时积分 $I$ 的值。
\end{enumerate}

\paragraph{四、应用题（9 分）}
在力场 $\vec F=(yz,zx,xy)$ 的作用下，质点由原点沿直线运动到椭球面
\[
  \frac{x^2}{a^2}+\frac{y^2}{b^2}+\frac{z^2}{c^2}=1
\]
上第一卦限的点 $M=(\xi,\eta,\zeta)$。问当 $\xi,\eta,\zeta$ 取何值时力场 $\vec F$ 做的功最大？最大值是多少？

\paragraph{五、证明题（共 2 小题，每小题 9 分，共 18 分）}
\begin{enumerate}
  \item 证明曲线
  \[
    \begin{cases}
      x^2+y^2+z^2-3x=0,\\
      2x-3y+5z-4=0
    \end{cases}
  \]
  在点 $(1,1,1)$ 处的法平面方程为 $16x+9y-z-24=0$。
  \item 证明级数
  \[
    \sum_{n=1}^\infty \sin\frac{\ln n^x}{n^3+x^2}
  \]
  在 $(-\infty,+\infty)$ 内一致收敛。
\end{enumerate}

\subsection{答案与解析}

\paragraph{一、填空题}
\begin{enumerate}
  \item $y=(x+C)\cos x$。
  \item $\sqrt2\,e$。
  \item $\dfrac{2}{15}$。
  \item
  \[
    \frac{a(3\sqrt3-1)}{8}
    +\frac{3a}{16}\ln\frac{3+2\sqrt3}{3}.
  \]
  \item $\dfrac23$。
\end{enumerate}

\paragraph{二、计算题}
\begin{enumerate}
  \item 通解为
  \[
    y=C_1\cos x+C_2\sin x+\frac32(x-1)e^x.
  \]
  \item 令
  \[
    w=xy+\frac{x}{y}.
  \]
  因 $z_x=(y+y^{-1})\cos w$，故
  \[
    z_{xy}
    =\left(1-\frac1{y^2}\right)\cos w
    -\left(y+\frac1y\right)\left(x-\frac{x}{y^2}\right)\sin w.
  \]
  \item 作变换
  \[
    u=x^2+y^2,\qquad v=y^2-x^2.
  \]
  第一象限内 $du\,dv=8xy\,dx\,dy$，区域变为 $9\le u\le16,\ 1\le v\le9$，故
  \[
    \iint_D xy\,dx\,dy
    =\frac18(16-9)(9-1)=7.
  \]
  \item 因 $\Omega$ 以原点为内点，而 $\vec F=\vec r/\|\vec r\|^3$ 在原点有孤立奇点，去掉原点附近小球并用高斯公式可得
  \[
    \iint_{\partial\Omega}\vec F\cdot\vec n\,dS=4\pi.
  \]
\end{enumerate}

\paragraph{三、解答题}
\begin{enumerate}
  \item
  \[
    \sum_{n=1}^\infty\frac1{4n(4n+2)}
    =\frac18\sum_{n=1}^\infty\frac1{n(2n+1)}
    =\frac{1-\ln2}{4}.
  \]
  \item 令 $\vec F=(y^2,z^2,x^2)$，则
  \[
    \nabla\times\vec F=(-2z,-2x,-2y).
  \]
  在平面 $x+y+z=1$ 上取与 $(1,1,1)$ 同向的单位法向量，则
  \[
    (\nabla\times\vec F)\cdot\vec n=-\frac2{\sqrt3}.
  \]
  三角形面积为 $\sqrt3/2$，由 Stokes 公式得所做的功为 $-1$。本题原卷有图，当前转写按题干给出的正向约定处理。
  \item 记
  \[
    P=\left[\sin x-\phi(x)\right]\frac{y}{x},\qquad Q=\phi(x).
  \]
  路径无关要求
  \[
    P_y=Q_x,\qquad \phi'(x)+\frac{\phi(x)}{x}=\frac{\sin x}{x}.
  \]
  因而 $(x\phi(x))'=\sin x$。由 $\phi(\pi)=1$ 得
  \[
    \phi(x)=\frac{\pi-1-\cos x}{x}.
  \]
  势函数可取 $U(x,y)=\phi(x)y$，故
  \[
    I=U(\pi,\pi)-U(1,0)=\pi.
  \]
\end{enumerate}

\paragraph{四、应用题}
沿直线从原点到 $M=(\xi,\eta,\zeta)$，取参数 $\vec r(t)=tM$，$0\le t\le1$。则
\[
  W=\int_0^1 3t^2\xi\eta\zeta\,dt=\xi\eta\zeta.
\]
在约束
\[
  \frac{\xi^2}{a^2}+\frac{\eta^2}{b^2}+\frac{\zeta^2}{c^2}=1
\]
下，令 $X=\xi/a,\ Y=\eta/b,\ Z=\zeta/c$，由对称性或拉格朗日乘数法知 $XYZ$ 在
\[
  X=Y=Z=\frac1{\sqrt3}
\]
处最大。故
\[
  M=\left(\frac a{\sqrt3},\frac b{\sqrt3},\frac c{\sqrt3}\right),\qquad
  W_{\max}=\frac{abc}{3\sqrt3}.
\]

\paragraph{五、证明题}
\begin{enumerate}
  \item 设
  \[
    F=x^2+y^2+z^2-3x,\qquad G=2x-3y+5z-4.
  \]
  在 $(1,1,1)$ 处，
  \[
    \nabla F=(-1,2,2),\qquad \nabla G=(2,-3,5).
  \]
  切向量可取
  \[
    \nabla F\times\nabla G=(16,9,-1).
  \]
  法平面垂直于切向量，故
  \[
    16(x-1)+9(y-1)-(z-1)=0,
  \]
  即 $16x+9y-z-24=0$。
  \item 因 $\ln n^x=x\ln n$，对任意 $x\in\mathbb R$ 有
  \[
    \left|\sin\frac{\ln n^x}{n^3+x^2}\right|
    \le \frac{|x|\ln n}{n^3+x^2}
    \le \frac{\ln n}{2n^{3/2}}.
  \]
  而 $\sum_{n=1}^\infty \ln n/n^{3/2}$ 收敛，故由 Weierstrass 判别法，原级数在 $(-\infty,+\infty)$ 上一致收敛。
\end{enumerate}
